\section{Diagram 3}\label{sec:d3}

Diagram~3 of \cite[Fig.~11]{AK} is the diagram of Roberts' continuum at infinity
\cite[\S7.2]{AK}.  With the numbering of \cite[\S5.2 and \S8.3]{AK}, where the clustering scheme
is $z:12\ldots34$, $w:14\ldots23$, it is the quadrilateral $1234$ with $z$-edges $14$, $23$ and
$w$-edges $12$, $34$; the diagonals and body $5$ carry no stroke.  Exchanging $z$ and $w$ gives
the same diagram with the bodies $2$ and $4$ exchanged, so after renumbering a Diagram-3 germ has
exactly this diagram.  Unlike \cite[proof of Thm.~6(iii)]{AK} we take no second singular limit.

\subsection{The crossed-cluster chart}
The equations of \eqref{eq:system} other than the gauge are invariant under
$(z,w,\delta)\mapsto(\mu z,\mu^{-1}w,\delta)$, $\mu\in\C^*$, and here we use the gauge
$z_{13}=w_{13}$, as \cite[p.~543]{AK} does for Roberts' continuum.  The chart has the parameter
$t$ and the coordinates $A,B,x,y,\hat z_5,C,D,u,v,\hat w_5$ and $D_{kl}$:
\begin{equation}\label{eq:d3-chart}
\begin{gathered}
z_1=\tfrac{A}{t},\quad z_2=\tfrac{A}{t}+tx,\quad z_3=\tfrac{B}{t},\quad z_4=\tfrac{B}{t}+ty,
\quad z_5=t^2\hat z_5,\\
w_1=\tfrac{C}{t},\quad w_2=\tfrac{D}{t},\quad w_3=\tfrac{D}{t}+tv,\quad w_4=\tfrac{C}{t}+tu,
\quad w_5=t^2\hat w_5,\\
\delta_{kl}=D_{kl}\ \ (kl=12,14,23,34),\qquad
\delta_{kl}=t\,D_{kl}\ \ (kl=13,24,15,25,35,45).
\end{gathered}
\end{equation}
Multiplying the $z$- and $w$-equations of the bodies $1,\dots,4$ and the gauge equation by $t$,
and the equations of body $5$ by $t^{-2}$, gives $21$ polynomial equations
$F=F^0+t^2F^2+\dots+t^6F^6=0$, with no term linear in $t$.  The \emph{boundary equations}
$F^0=0$ say that the following expressions vanish:
\begin{equation}\label{eq:d3-boundary}
\begin{gathered}
R^z_1=A-m_4D_{14}^3(A-B),\qquad R^z_4=B-m_1D_{14}^3(B-A),\\
R^z_2=A-m_3D_{23}^3(A-B),\qquad R^z_3=B-m_2D_{23}^3(B-A),\\
R^w_1=C-m_2D_{12}^3(C-D),\qquad R^w_2=D-m_1D_{12}^3(D-C),\\
R^w_3=D-m_4D_{34}^3(D-C),\qquad R^w_4=C-m_3D_{34}^3(C-D),\\
\hat z_5+(m_1D_{15}^3+m_2D_{25}^3)A+(m_3D_{35}^3+m_4D_{45}^3)B,\\
\hat w_5+(m_1D_{15}^3+m_4D_{45}^3)C+(m_2D_{25}^3+m_3D_{35}^3)D,\\
D_{12}^2x(D-C)-1,\quad D_{34}^2y(C-D)-1,\quad D_{14}^2u(B-A)-1,\\
D_{23}^2v(B-A)-1,\quad D_{13}^2(B-A)(D-C)-1,\quad D_{24}^2(B-A)(C-D)-1,\\
D_{15}^2AC-1,\quad D_{25}^2AD-1,\quad D_{35}^2BD-1,\quad D_{45}^2BC-1,\quad (B-A)-(D-C);
\end{gathered}
\end{equation}
$R^z_k$, $R^w_k$ come from the equations of body $k$.  We work on the slice $B=A+1$; for a germ
we take $t=1/z_{13}$ (Proposition~\ref{prop:d3-cover}).

\subsection{Coverage}
We write $s$ for the parameter of a germ.

\begin{proposition}[Coverage]\label{prop:d3-cover}
Let $\gamma$ be a Diagram-3 germ, renormalized as in Lemma~\ref{lem:germ}.
\begin{romanlist}
  \item Bodies $1,\dots,4$ are circled in both colours, body $5$ in neither.
  \item The sides are maximal edges: $r_{12},r_{14},r_{23},r_{34}\approx1$,
  $z_{14},z_{23},w_{12},w_{34}\approx\epsilon^{-2}$ and $w_{14},w_{23},z_{12},z_{34}\approx\epsilon^2$.
  \item $z_{13},w_{13},z_{24},w_{24},z_{k5},w_{k5}\approx\epsilon^{-2}$ for $k\le4$, and
  $z_5,w_5\preccurlyeq\epsilon^4$.
  \item Renormalize instead by $\lambda=(w_{13}/z_{13})^{1/2}$, so that $z_{13}=w_{13}$, and put
  $t=1/z_{13}\approx\epsilon^2$.  Then $B-A=1$ and every chart coordinate converges as $s\to0$; the
  limit $X_0$ satisfies \eqref{eq:d3-boundary}, and $A,B,C,D,x,y,u,v$ and all $D_{kl}$ are
  nonzero at $X_0$.
\end{romanlist}
\end{proposition}

\begin{proof}
By Lemma~\ref{lem:germ} the relations of \cite[\S3]{AK} hold between orders along $\gamma$, and
Fact~\ref{fact:rules} applies.  (i) Each side is the only stroke of its colour at its two ends, so
Rule~1a \cite[pp.~545--546]{AK} circles $1,\dots,4$ in both colours; a circle needs a stroke of
its colour, so $5$ has none.
(ii) The isolated component $\{1,4\}$ of the $z$-diagram has $z$-circles, so its stroke $14$ is
maximal by Rule~1e \cite[p.~546]{AK}; likewise $23$, and $12$, $34$ in the $w$-diagram.  The right
equality case of Estimate~2 \cite[p.~548]{AK} gives the orders.
(iii) By (ii), $z_{13}=z_{14}-z_{34}$, $w_{13}=w_{12}+w_{23}$, $z_{24}=z_{23}+z_{34}$ and
$w_{24}=w_{14}-w_{12}$ are $\approx\epsilon^{-2}$; by (i), $z_k\approx Z=\epsilon^{-2}\succ z_5$ for
$k\le4$, and likewise for $w$.  The equation
$z_5(1-\sum_{k\le4}m_k\delta_{k5}^3)=-\sum_{k\le4}m_k\delta_{k5}^3z_k$ of body $5$, with
$\delta_{k5}\approx\epsilon^2$, gives $z_5\preccurlyeq\epsilon^4$; likewise for $w_5$.
(iv) By (iii), $\lambda$ differs from the factor of Lemma~\ref{lem:germ} by a Puiseux series of
order $0$, so (i)--(iii) persist.  Now $B-A=tz_{13}=1$, and $A=z_1/z_{13}$, $B=z_3/z_{13}$,
$C=w_1/w_{13}$, $D=w_2/w_{13}$, $x=z_{12}z_{13}$, $y=z_{34}z_{13}$, $u=w_{14}w_{13}$,
$v=w_{23}w_{13}$ and the $D_{kl}$ ($=\delta_{kl}$ for the sides, $\delta_{kl}z_{13}$ otherwise) are
$\approx1$, while $\hat z_5=z_5z_{13}^2$, $\hat w_5=w_5w_{13}^2\preccurlyeq1$.  So all chart
coordinates converge, and $F=0$ passes to the limit.
\end{proof}

By Proposition~\ref{prop:d3-cover}(ii) and~(iii), every Diagram-3 germ has an unbounded
four-product: $p_{12}=r_{12}^2r_{34}^2r_{35}^2r_{45}^2\approx\epsilon^{-8}$.

\subsection{The first condition}

\begin{proposition}[First condition]\label{prop:d3-first}
At every point of $\{t=0\}$ satisfying \eqref{eq:d3-boundary}, $m_1m_3=m_2m_4$.  Hence every
Diagram-3 germ satisfies (29) of \cite[\S8.3]{AK}.
\end{proposition}

\begin{proof}
Since $m_1R^z_1+m_4R^z_4=m_1A+m_4B$ and $m_2R^z_2+m_3R^z_3=m_2A+m_3B$ identically,
\begin{equation}\label{eq:d3-syzygy}
(B-A)(m_1m_3-m_2m_4)=(m_1+m_4)(m_2R^z_2+m_3R^z_3)-(m_2+m_3)(m_1R^z_1+m_4R^z_4),
\end{equation}
and likewise
$(D-C)(m_1m_3-m_2m_4)=(m_1+m_2)(m_3R^w_3+m_4R^w_4)-(m_3+m_4)(m_1R^w_1+m_2R^w_2)$.  With
$T_{13}=D_{13}^2(B-A)(D-C)-1$ (the distance equation of $13$) and $g=(B-A)-(D-C)$ (the gauge),
$D_{13}^2(B-A)^2-1=T_{13}+D_{13}^2(B-A)g$.  At a boundary point $T_{13}=g=0$, so
$D_{13}^2(B-A)^2=1$, $B-A\ne0$, and \eqref{eq:d3-syzygy} gives $m_1m_3=m_2m_4$.  The second
statement follows from Proposition~\ref{prop:d3-cover}(iv).
\end{proof}

These identities are the chart form of \cite[\S5.2]{AK}, where Rule~1c makes $m_1z_1+m_4z_4$ and
$m_2z_2+m_3z_3$ non-maximal while $z_1\sim z_2$ and $z_3\sim z_4$.

\subsection{Labels and the analytic graph}
From now on $m\in\Rp^5$ satisfies (29), $M=m_1+m_2+m_3+m_4$, and $X$ denotes the $19$ chart
coordinates other than $B=A+1$.  Let $F_{z_k},F_{w_k}$ be the rescaled equations of body $k$, let
$\Lambda$ be the $19\times21$ matrix that forms from $F$ the combinations $m_1F_{z_1}+m_4F_{z_4}$,
$F_{z_4}-F_{z_1}$, $F_{z_3}-F_{z_2}$, $m_1F_{w_1}+m_2F_{w_2}$, $F_{w_2}-F_{w_1}$,
$F_{w_3}-F_{w_4}$ and the other $13$ equations, and put
$\ell_zF=F_{z_1}-F_{z_2}-\frac{m_3}{m_2}F_{z_3}+\frac{m_4}{m_1}F_{z_4}$,
$\ell_wF=F_{w_1}+\frac{m_2}{m_1}F_{w_2}-\frac{m_2}{m_1}F_{w_3}-F_{w_4}$.

\begin{lemma}[computer-assisted]\label{lem:d3-fibre}
\begin{alphalist}
  \item On the slice, \eqref{eq:d3-boundary} has exactly $3^4\cdot2^6=5184$ solutions, the
  \emph{labels}: $D_{12},D_{34},D_{14},D_{23}$ are arbitrary cube roots of $\frac1{m_1+m_2}$,
  $\frac1{m_3+m_4}$, $\frac1{m_1+m_4}$, $\frac1{m_2+m_3}$; $D_{13},D_{24},D_{15},D_{25},D_{35},
  D_{45}$ are arbitrary square roots of $1$, $-1$, $\frac M{m_3}$, $-\frac M{m_4}$, $\frac M{m_1}$,
  $-\frac M{m_2}$; $A=-\frac{m_4}{m_1+m_4}$, $C=-\frac{m_2}{m_1+m_2}$, $D=\frac{m_1}{m_1+m_2}$,
  $x=(m_1+m_2)D_{12}$, $y=-(m_3+m_4)D_{34}$, $u=(m_1+m_4)D_{14}$, $v=(m_2+m_3)D_{23}$, and
  $\hat z_5,\hat w_5$ follow from the equations of body $5$.  The mass $m_5$ does not enter.
  \item At every label
  \[
  \det\partial_X(\Lambda F^0)=-\frac{5184\,m_1^2m_2^3m_4^3\,(D_{12}D_{14}D_{23}D_{34})^4\,
  D_{13}D_{24}D_{15}D_{25}D_{35}D_{45}}{(m_1+m_2)(m_1+m_4)}\neq0 .
  \]
\end{alphalist}
\end{lemma}

The proof is an exact computation over the field $\mathbb Q(m_1,m_2,m_4,m_5)$, with
$m_3=m_2m_4/m_1$.  On the slice, the $19$ entries of $\Lambda F^0$ and the relations listed in (a)
are combinations of each other with coefficients that have no poles at positive masses, and the
determinant in (b) is computed by exact elimination (Appendix~\ref{app:verification}).
Identically in $X$, $\ell_zF^0=-B(m_1m_3-m_2m_4)/(m_1m_2)$ and
$\ell_wF^0=-(D-C)D_{34}^3(m_1m_3-m_2m_4)/m_1$, so $\ell_z,\ell_w$ annihilate
$J_0=\partial_XF^0(X_0)$; as $\Lambda J_0$ is invertible, $J_0$ has rank $19$, $\ell_z,\ell_w$ span
its left kernel, and the $21$ rows of $\Lambda,\ell_z,\ell_w$ are independent.

\begin{lemma}\label{lem:d3-graph}
Let $X_0$ be a label.  There is a unique analytic $t\mapsto X(t)$ near $0$ with $X(0)=X_0$ and
$\Lambda F(t,X(t))=0$; put $G(t)=(\ell_zF(t,X(t)),\ell_wF(t,X(t)))=\sum_{k\ge1}G_kt^k$.  If the
chart coordinates of a Diagram-3 germ converge to $X_0$, then along the germ they are $X(t)$, and
$G\equiv0$.  Every analytic $Y(t)$ with $Y(0)=X_0$ and $F(t,Y(t))\equiv0$ equals $X(t)$.
\end{lemma}

\begin{proof}
The first statement is the analytic implicit function theorem (Lemma~\ref{lem:d3-fibre}(b)).
Along the germ, after a reparametrization $s\mapsto s^N$, the chart coordinates $Y(s)$ and
$t(s)=1/z_{13}(s)$ are convergent power series (Lemma~\ref{lem:germ},
Proposition~\ref{prop:d3-cover}) with $Y(0)=X_0$, $t(0)=0$, $t\not\equiv0$.  Since
$\Lambda F(t(s),Y(s))=0$, uniqueness gives $Y(s)=X(t(s))$.  So $F(t,X(t))$ vanishes on the image
of $t(s)$, a neighbourhood of $0$; hence $F(t,X(t))\equiv0$ and $G\equiv0$.  The last statement is
again uniqueness.
\end{proof}

So a label carries a Diagram-3 germ, ramified or not, only if $G\equiv0$ there.  As $F$ has no
term linear in $t$, $X(t)=X_0+X_2t^2+X_3t^3+\cdots$, and as $\ell J_0=0$ ($\ell=(\ell_z,\ell_w)$),
$G_k$ does not involve $X_k$: at $X_0$, $G_2=\ell F^2$, $G_3=\ell F^3$,
$G_4=\ell(F^4+\partial F^2\cdot X_2+\tfrac12\partial^2F^0[X_2,X_2])$ and
$G_5=\ell(F^5+\partial F^3\cdot X_2+\partial F^2\cdot X_3+\partial^2F^0[X_2,X_3])$.

\subsection{Orders three and four}

\begin{lemma}[computer-assisted]\label{lem:d3-order3}
At every label $G_2=0$ and
$G_3=\bigl(\frac{m_1+m_2}{m_1^2m_2}\Omega,\frac{m_1+m_4}{m_1^2m_4}\Omega\bigr)$, where
$\Omega=\Omega_{\mathrm{re}}+\Omega_{\mathrm{im}}$,
\[
\begin{aligned}
\Omega_{\mathrm{re}}&=m_1m_2m_4D_{13}+m_5m_1^2D_{15}+m_5m_2m_4D_{35},\\
\Omega_{\mathrm{im}}&=m_1m_2m_4D_{24}+m_5m_1m_2D_{25}+m_5m_1m_4D_{45}.
\end{aligned}
\]
\end{lemma}

The computation is exact, over the field $\mathbb Q(m_1,m_2,m_4,m_5)$ with $m_3=m_2m_4/m_1$:
$\ell F^2$ and $\ell F^3$ are evaluated at the labels and reduced modulo the cubic and quadratic
equations satisfied by the $D_{kl}$ (Appendix~\ref{app:verification}).

\begin{corollary}\label{cor:d3-order3}
If $G_3=0$ at some label, then $N_{\mathrm{re}}(m)=N_{\mathrm{im}}(m)=0$, where $Q=m_5^2M$ and
\begin{align*}
N_{\mathrm{re}}&=\bigl((m_1m_2m_4)^3-Q(m_1^6+m_2^3m_4^3)\bigr)^2-4Q^2m_1^6m_2^3m_4^3,\\
N_{\mathrm{im}}&=\bigl((m_1m_2m_4)^3-Qm_1^3(m_2^3+m_4^3)\bigr)^2-4Q^2m_1^6m_2^3m_4^3 .
\end{align*}
As $N_{\mathrm{re}}-N_{\mathrm{im}}+QS=0$ with
$S=(m_1^3-m_2^3)(m_1^3-m_4^3)\bigl(2(m_1m_2m_4)^3-Q(m_1^3+m_2^3)(m_1^3+m_4^3)\bigr)$, also
$S(m)=0$.
\end{corollary}

\begin{proof}
At positive masses $D_{13},D_{15},D_{35}$ are real and $D_{24},D_{25},D_{45}$ purely imaginary
(Lemma~\ref{lem:d3-fibre}), so $\Omega_{\mathrm{re}}=\Omega_{\mathrm{im}}=0$.  If
$\alpha+\beta+\gamma=0$ then $(\alpha^2-\beta^2-\gamma^2)^2-4\beta^2\gamma^2=\prod_\pm(\alpha+\beta
\pm\gamma)\prod_\pm(\alpha-\beta\pm\gamma)=0$.  Inserting the squares $(m_1m_2m_4)^2$,
$m_5^2m_1^4M/m_3$, $m_5^2m_2^2m_4^2M/m_1$ of the terms of $\Omega_{\mathrm{re}}$, using
$m_3=m_2m_4/m_1$ and multiplying by $(m_1m_2m_4)^2$ gives $N_{\mathrm{re}}=0$; the terms of
$\Omega_{\mathrm{im}}$ are $i$ times reals with squares $(m_1m_2m_4)^2$, $m_5^2m_1^2m_2^2M/m_4$,
$m_5^2m_1^2m_4^2M/m_2$, which gives $N_{\mathrm{im}}=0$.  Finally $Q>0$.
\end{proof}

\begin{lemma}[computer-assisted]\label{lem:d3-order4}
At every label $G_4=-\frac{\Pi}{m_1^2}\bigl(m_4(m_1+m_2),m_2(m_1+m_4)\bigr)$ with
$\Pi=(m_1D_{12}-m_4D_{34})(m_1D_{14}-m_2D_{23})$.
\end{lemma}

As for Lemma~\ref{lem:d3-order3}, the computation is exact and reduces modulo the equations
satisfied by the $D_{kl}$; it solves for $X_2$ and evaluates the formula for $G_4$ above
(Appendix~\ref{app:verification}).  At every label
$(m_1D_{12})^3-(m_4D_{34})^3=m_1(m_1+m_4)(m_1-m_4)/(m_1+m_2)$ and
$(m_1D_{14})^3-(m_2D_{23})^3=m_1(m_1+m_2)(m_1-m_2)/(m_1+m_4)$, so $\Pi=0$ forces $m_1=m_4$ or
$m_1=m_2$; with (29), $m=(p,q,q,p,m_5)$ or $(p,p,q,q,m_5)$.

\begin{lemma}[computer-assisted]\label{lem:d3-curves}
Let $m\in\Rp^5$ satisfy (29), $N_{\mathrm{re}}(m)=N_{\mathrm{im}}(m)=0$ and $m_1=m_2$ (resp.\
$m_1=m_4$).  Write $m_1=a^2$, $m_3=b^2$ (resp.\ $m_2=b^2$) with $a,b>0$, and
$c=(2(a^2+b^2))^{1/2}$.  Then
\[
m_5=\frac{a^3b^3}{c\,|a^3+b^3|}\qquad\text{or}\qquad m_5=\frac{a^3b^3}{c\,|a^3-b^3|}\quad(a\ne b).
\]
Conversely $N_{\mathrm{re}}$, $N_{\mathrm{im}}$ and $S$ vanish on these four families, called
$E^\pm_{12}$ (if $m_1=m_2$) and $E^\pm_{14}$ (if $m_1=m_4$).  They meet only at $a=b$, i.e.\ at
$m$ proportional to $(1,1,1,1,\frac14)$, which lies on $E^+_{12}\cap E^+_{14}$.
\end{lemma}

On these mass patterns $N_{\mathrm{re}}=a^{12}\prod_\pm\bigl(a^6b^6-m_5^2c^2(a^3\pm b^3)^2\bigr)$,
a polynomial identity in $a,b,m_5$.  The converse is checked by exact substitution, with $a,b,c$
given by a rational parametrization of the conic $c^2=2(a^2+b^2)$
(Appendix~\ref{app:verification}).

\begin{lemma}[computer-assisted]\label{lem:d3-census}
Let $m$ lie on one of the four families.
\begin{alphalist}
  \item If $a\ne b$, exactly $108$ labels have $G_3=G_4=0$: two of the eight sign choices for
  $(D_{13},D_{15},D_{35})$ give $\Omega_{\mathrm{re}}=0$, two of those for $(D_{24},D_{25},D_{45})$
  give $\Omega_{\mathrm{im}}=0$, and $\Pi=0$ exactly when $D_{14}=D_{23}$ (on $E^\pm_{12}$) or
  $D_{12}=D_{34}$ (on $E^\pm_{14}$), for $27$ of the $81$ choices of cube roots.
  \item If $a=b$, exactly $180$ labels have $G_3=G_4=0$: those with $D_{15}=D_{35}=-2D_{13}$,
  $D_{25}=D_{45}=-2D_{24}$, and $D_{12}=D_{34}$ or $D_{14}=D_{23}$ ($45$ choices of cube roots).
\end{alphalist}
\end{lemma}

An exact computation, with $a,b,c$ as indeterminates, finds the sign choices for which the
relevant linear forms in $a^3,b^3$ vanish identically (Appendix~\ref{app:verification}).  The
count holds at every point: on $E^\pm_{12}$, with $D_{13}=\sigma$, $D_{15}=s_1c/b$,
$D_{35}=s_2c/a$ and $\nu=\pm1$ the sign of $a^3\pm b^3$, $\Omega_{\mathrm{re}}$ is
$a^4b^2/|a^3\pm b^3|$ times $(\sigma\nu+s_1)a^3+(\pm\sigma\nu+s_2)b^3$, whose coefficients lie in
$\{0,2\sigma\nu\}$
and $\{0,\pm2\sigma\nu\}$; so for $a,b>0$ (and $a\ne b$ on $E^-_{12}$) it vanishes only if both
coefficients do.  The other cases are alike.

\subsection{Orders five and six}

\begin{lemma}[computer-assisted]\label{lem:d3-order5}
Let $m$ lie on one of the four families with $a\ne b$.  Then $G_5\neq0$ at each of the $108$
labels of Lemma~\ref{lem:d3-census}(a).  At $a=b$, $G_5=0$ at all $180$ labels of
Lemma~\ref{lem:d3-census}(b).
\end{lemma}

We compute $G_5$ in exact arithmetic over the field $\mathbb Q(i,a,b,c)$, $c^2=2(a^2+b^2)$, modulo
the cubic equations satisfied by $D_{12},D_{34},D_{14}$, and the product over the nine pairs below
as an iterated resultant (Appendix~\ref{app:verification}).  On $E^+_{12}$ the labels have
$D_{14}=D_{23}$, and the $z$-component of $G_5$ is a nonzero factor
times $(a-b)\,L(D_{12}/D_{14},D_{34}/D_{14})$, where
$L(\xi_1,\xi_2)=a^2F_4\xi_1-a^{-2}b^4F_4\xi_2-b^2(a^2-b^2)F_2$,
$F_4=a^4+a^3b+a^2b^2+ab^3+b^4$ and $F_2=a^2+ab+b^2$.  Up to a nonzero monomial factor and sign,
the product of $L$ over the nine pairs with $\xi_1^3=(a^2+b^2)/(2a^2)$, $\xi_2^3=(a^2+b^2)/(2b^2)$
is $(a-b)^3(a+b)^3P_{66}(a,b)$, where $P_{66}$ is homogeneous of degree $66$ with $67$ positive
coefficients; so $G_5\ne0$ if $a\ne b$.
On $E^-_{12}$ the same computation gives $P_{66}(a,-b)$, and an exact Sturm count shows that
$P_{66}(1,-\rho)$ has no root $\rho>0$.  Exchanging the bodies $2$ and $4$ together with $z$ and
$w$ maps the chart at $m$ onto the chart at $(m_1,m_4,m_3,m_2,m_5)$, $E^\pm_{12}$ onto
$E^\pm_{14}$ and the $108$ labels onto the $108$ labels; the relabelling $(13)(24)$ exchanges $a$
and $b$ and covers the minus families with $a<b$.  At $a=b$ both components of $G_5$ vanish at
the labels of $E^+_{12}$, hence, by the exchange, at all $180$ labels.

\begin{lemma}[computer-assisted]\label{lem:d3-order6}
Let $m=(\frac12,\frac12,\frac12,\frac12,\frac18)$.  Of the $180$ labels of
Lemma~\ref{lem:d3-census}(b), the $144$ labels with $D_{12}\ne D_{34}$ or $D_{14}\ne D_{23}$ have
$G_6\neq0$.
\end{lemma}

The proof is a computation of the jets modulo two primes (Appendix~\ref{app:verification}).  All
data lie in $R=\mathbb Z[\frac12,i,\omega]$, $\omega=e^{2\pi i/3}$; let $\mathfrak p$ be a prime of
$R$ above $p\in\{13,37\}$, so $R/\mathfrak p=\mathbb F_p$ as $p\equiv1\pmod{12}$.  If a germ passed
through a label, then $F(t,X(t))\equiv0$ (Lemma~\ref{lem:d3-graph}); each $X_k$ would be the
unique solution of the order-$k$ equations, given by any $19$ rows with invertible minor, and all
$21$ residuals would vanish.  We solve with $19$ rows whose minor is a unit modulo $\mathfrak p$
and find a nonzero residual at order $6$ for the $144$ labels.  Another choice of $\mathfrak p$
acts on the labels by $\mathrm{Gal}(\mathbb Q(i,\omega)/\mathbb Q)$ ($\omega\mapsto\omega^2$,
$i\mapsto-i$), which preserves the set of the $144$ labels; both primes give the same set.

\begin{lemma}[computer-assisted]\label{lem:d3-twisted}
Let $m=(\frac12,\frac12,\frac12,\frac12,\frac18)$, $\alpha,\beta\in\{1,\omega,\omega^2\}$ and
$\varepsilon,\eta\in\{\pm1\}$.  Let $(P,Q,Y)$ be the analytic solution near $\tau=0$ of
\[
(1-P)(1+Q)=\alpha\tau^2,\qquad (1+P)(1-Q)=\beta\tau^2,\qquad 4PQY^2+1=0
\]
with $(P,Q,Y)(0)=(1,1,\eta i/2)$; the Jacobian there is $\pm16i$.  Then
\[
z=\tfrac1\tau(1,P,-1,-P,0),\qquad w=\tfrac1\tau(1,-Q,-1,Q,0),
\]
\[
\delta_{12}=\delta_{34}=\alpha,\ \ \delta_{14}=\delta_{23}=\beta,\ \
\delta_{13}=\tfrac{\varepsilon\tau}2,\ \ \delta_{15}=\delta_{35}=-\varepsilon\tau,\ \
\delta_{24}=\tau Y,\ \ \delta_{25}=\delta_{45}=-2\tau Y
\]
solves \eqref{eq:system} with the gauge $z_{13}=w_{13}$.  In the chart, with $t=-\tau/2$, it is an
analytic solution of $F=0$ through the label with $D_{12}=D_{34}=\alpha$, $D_{14}=D_{23}=\beta$,
$D_{13}=-\varepsilon$ and $D_{24}=-\eta i$.  These $36$ labels are the labels of
Lemma~\ref{lem:d3-census}(b) with $D_{12}=D_{34}$ and $D_{14}=D_{23}$.  Along these germs
\begin{equation}\label{eq:d3-identities}
\delta_{15}+2\delta_{13}=\delta_{35}+2\delta_{13}=\delta_{25}+2\delta_{24}=\delta_{45}+2\delta_{24}=0 .
\end{equation}
\end{lemma}

We substitute into \eqref{eq:system} in exact arithmetic, with $\alpha,\beta,\varepsilon,\tau,P,Q,Y$
as indeterminates: the coordinate equations vanish identically once $\alpha^3=\beta^3=1$, the
distance equations reduce to the three displayed ones, and \eqref{eq:d3-identities} holds without
using them (Appendix~\ref{app:verification}).  The chart coordinates are analytic in $t$ because
$P,Q,Y$ are analytic in $\tau^2$ with $P(0)=Q(0)=1$.

\subsection{The exclusion}

\begin{theorem}\label{thm:d3}
Let $m\in\Rp^5$, let $\Gamma\subset A_m$ be an irreducible curve, and let $\gamma$ be a Diagram-3
germ on $\Gamma$, with the bodies numbered as in Diagram~3.  Then $m_1=m_2=m_3=m_4=4m_5$, and
\eqref{eq:d3-identities} holds identically on $\Gamma$.  In particular $\Gamma$ contains no
positive point, and Theorem~\ref{thm:exclusion}(c) holds.
\end{theorem}

\begin{proof}
Exchanging $z$ and $w$ and renumbering if necessary, we may assume that $\gamma$ has the diagram
described at the beginning of this section; the conclusions are invariant under these
operations.  By Proposition~\ref{prop:d3-cover} the chart coordinates of $\gamma$ converge to a
solution $X_0$ of \eqref{eq:d3-boundary} with $B-A=1$.  By Proposition~\ref{prop:d3-first},
(29) holds, so $X_0$ is a label (Lemma~\ref{lem:d3-fibre}), and $G\equiv0$ at $X_0$
(Lemma~\ref{lem:d3-graph}).  By Lemma~\ref{lem:d3-order3} and Corollary~\ref{cor:d3-order3},
$N_{\mathrm{re}}(m)=N_{\mathrm{im}}(m)=0$.  By Lemma~\ref{lem:d3-order4}, $m_1=m_4$ or $m_1=m_2$,
so $m$ lies on one of the families of Lemma~\ref{lem:d3-curves}.  By
Lemmas~\ref{lem:d3-census}(a) and~\ref{lem:d3-order5}, $a=b$, that is $m_1=m_2=m_3=m_4=4m_5$.

The system \eqref{eq:system} is covariant under
$(m,z,w,\delta)\mapsto(\kappa^3m,\kappa z,\kappa w,\kappa^{-1}\delta)$, $\kappa>0$, which maps
Diagram-3 germs to Diagram-3 germs and preserves \eqref{eq:d3-identities}.  So we may take
$m=(\frac12,\frac12,\frac12,\frac12,\frac18)$.
By Lemmas~\ref{lem:d3-census}(b) and~\ref{lem:d3-order6}, $D_{12}=D_{34}$ and $D_{14}=D_{23}$ at
$X_0$.  By Lemmas~\ref{lem:d3-twisted} and~\ref{lem:d3-graph}, the chart coordinates of $\gamma$
are those of the twisted Roberts germ through $X_0$.  The renormalizations do not change the
$\delta_{kl}$, so \eqref{eq:d3-identities} holds along $\gamma$.  Each expression in
\eqref{eq:d3-identities} is a regular function on the irreducible curve $\Gamma$ that vanishes at
the infinitely many points of $\gamma$, so it vanishes on $\Gamma$.  At a positive point
$\delta_{15}+2\delta_{13}>0$.
\end{proof}

So every Diagram-3 germ has the masses $(4,4,4,4,1)$ up to scale and is, up to the symmetries
used above, one of the $36$ twisted Roberts germs, which do exist.  The next remark shows that
the theorem cannot be weakened to curves through a real point.

\begin{remark}[The modified Roberts continuum]\label{rem:d3-roberts}
Let $m=(\frac12,\frac12,\frac12,\frac12,\frac18)$, proportional to $(4,4,4,4,1)$
\cite[Rem.~8, p.~583]{AK}.  The curve
\[
z=(a,b,-a,-b,0),\quad w=(a,-b,-a,b,0),\quad b^2=a^2-1,
\]
\[
\delta_{12}=\delta_{23}=\delta_{34}=\delta_{14}=1,\ \ \delta_{13}=\tfrac1{2a},\ \
\delta_{24}=\tfrac{i}{2b},\ \ \delta_{15}=\delta_{35}=-\tfrac1a,\ \ \delta_{25}=\delta_{45}=-\tfrac ib
\]
is the modified Roberts continuum \cite[p.~543]{AK}.  It solves \eqref{eq:system} with the gauge
$z_{13}=w_{13}$, $U=I=1$ on it, and \eqref{eq:d3-identities} holds, so it has no positive point.
Since $w_{12}/z_{12}=(a+b)^2$ on it, the rotation by $\mu=a+b$ gives an irreducible curve
$\Gamma_R\subset A_m$.  For $0<a<1$ and $b\in i\R$, $|\mu|=1$ and the points are real, with
$\delta_{13}>0>\delta_{15}=\delta_{35}$.
As $a\to\infty$ with $b\sim a$, the curve is the twisted Roberts germ of
Lemma~\ref{lem:d3-twisted} with $(\alpha,\beta,\varepsilon,\eta)=(1,1,1,1)$, $\tau=1/a$ and
$P=Q=b/a$.  Its strokes are $z_{14}$, $z_{23}$, $w_{12}$, $w_{34}$, so it is a Diagram-3 germ
(Roberts' continuum at infinity \cite[\S7.2]{AK}); the sides keep length $1$,
$r_{13},r_{24},r_{k5}\to\infty$, $p_{12}=a^2(1-a^2)$ is unbounded, and (29) holds.
So $\Gamma_R$ passes through real points, carries the positive constant $U=I=1$, and has a
Diagram-3 germ with an unbounded four-product.  An exclusion at the level of curves through a
real point is therefore false for Diagram~3.  Theorem~\ref{thm:d3} uses the positivity of points
only through the identity $\delta_{15}+2\delta_{13}=0$ on all of $\Gamma$, as
Remark~\ref{rem:levels} requires.
\end{remark}
